% ---------------------------------------------------------------------
%  §2. Introducción
% ---------------------------------------------------------------------
\section{Introducción}

El uso de redes sociales forma parte de la vida cotidiana de los estudiantes, pero su impacto puede no ser
igual para todos. Hablar de “redes sociales” en general esconde que no todas se usan igual, y el contexto de
cada estudiante —el país en que vive, su género o su situación personal— puede influir en sus hábitos
digitales y en sus consecuencias. Identificar qué grupos y qué patrones de uso se asocian con peor
bienestar permite formular recomendaciones más específicas.

El trabajo utiliza el conjunto de datos \textit{Students Social Media Addiction}, publicado en Kaggle
\cite{kaggle}, con información de estudiantes secundarios, universitarios y de posgrado sobre su uso diario
de redes, la plataforma que más utilizan, sus horas de sueño, un puntaje de salud mental, su percepción
sobre si las redes afectan su rendimiento académico y su situación sentimental.

\begin{hallazgo}
\textbf{Pregunta principal:} \textit{¿cómo impacta el uso de redes sociales en los estudiantes y es ese
impacto igual para todos?} Se descompone en cuatro preguntas:
\begin{enumerate}[label=\arabic*.,itemsep=0pt,topsep=2pt]
  \item ¿El impacto es igual para todos los estudiantes según el nivel de desarrollo de su país?
  \item ¿El uso de redes sociales impacta en el sueño y en el rendimiento académico?
  \item ¿Hay redes sociales que impactan más que otras y existe una brecha de género en ese impacto?
  \item ¿El uso de redes sociales afecta las relaciones sentimentales de los estudiantes?
\end{enumerate}
\end{hallazgo}

En todo el informe, un puntaje de salud mental más alto representa un \textbf{mejor} estado de salud mental.
